%==============================================================================
\chapter{Caracterización de instancias y conjunto de datos}
\label{cap:caracterizacion}
%==============================================================================
% Objetivo específico 2. Actividades A5 a A10.

Para llevar a la práctica los modelos de MiniZinc y la formulación del
\textit{schedule} definidos en el capítulo anterior, es necesario
reunir un conjunto de instancias de prueba, extraer sus descriptores y medir
en ellas el comportamiento real de cada configuración. Ese proceso abarca la
conversión de las instancias (Sección~\ref{sec:instancias}), el cálculo de sus
cotas (Sección~\ref{sec:cotas_instancias}), la extracción de características
(Sección~\ref{sec:features}) y las campañas de ejecución de los
\textit{solvers} (Sección~\ref{sec:ejecucion_solvers}), cuyos registros se
consolidan en la base de conocimiento (Sección~\ref{sec:preprocesamiento}) que
alimenta las estrategias de selección del Capítulo~\ref{cap:portafolio}.

%------------------------------------------------------------------------------
\section{Conjunto de instancias}
\label{sec:instancias}
%------------------------------------------------------------------------------

\subsection{Familias de instancias}

Se trabajó con 242 instancias clásicas del JSSP, distribuidas en ocho familias
ampliamente utilizadas como \textit{benchmark} en la literatura. Varias de ellas
se distribuyen a través de la OR-Library \citep{beasley1990}, y todas están
reunidas, junto con sus cotas y mejores soluciones conocidas, en el repositorio
de \citet{weise2019jssp}. En la Tabla~\ref{tab:familias} se detallan el origen,
el número de instancias y las dimensiones de cada familia.

\begin{table}[H]
  \centering
  \small
  \singlespacing
  \caption{Familias de instancias utilizadas. El tamaño se expresa como
  $n \times m$ (trabajos $\times$ máquinas).}
  \label{tab:familias}
  \begin{tabularx}{\textwidth}{>{\ttfamily}l >{\raggedright\arraybackslash}X
                               c >{\raggedright\arraybackslash}p{0.33\textwidth}}
    \toprule
    \normalfont\textbf{Familia} & \textbf{Origen} & \textbf{Instancias} & \textbf{Tamaños} \\
    \midrule
    ft  & \citet{fisher1963}         & 3  & $6\times6$, $10\times10$, $20\times5$ \\
    la  & \citet{lawrence1984}       & 40 & de $10\times5$ a $30\times10$ \\
    abz & \citet{adams1988}          & 5  & $10\times10$, $20\times15$ \\
    orb & \citet{applegate1991}      & 10 & $10\times10$ \\
    swv & \citet{storer1992}         & 20 & $20\times10$, $20\times15$, $50\times10$ \\
    yn  & \citet{yamada1992}         & 4  & $20\times20$ \\
    ta  & \citet{taillard1993}       & 80 & de $15\times15$ a $100\times20$ \\
    dmu & \citet{demirkol1998}       & 80 & de $20\times15$ a $50\times20$ \\
    \midrule
    \normalfont\textbf{Total} & & 242 & de $6\times6$ a $100\times20$ \\
    \bottomrule
  \end{tabularx}
  \par\vspace{2pt}{\footnotesize El origen de las familias \texttt{ft}, \texttt{la},
  \texttt{abz}, \texttt{orb}, \texttt{swv} e \texttt{yn} es el que indica la
  OR-Library \citep{orlibjobshop}.}
\end{table}

El conjunto cubre un rango amplio de dificultad. Las instancias pequeñas de las
familias \texttt{ft}, \texttt{la} y \texttt{orb} se resuelven a optimalidad en
segundos con métodos exactos modernos, mientras que muchas instancias de
\texttt{ta}, \texttt{dmu}, \texttt{swv} e \texttt{yn} siguen abiertas, ya que
su óptimo no se conoce y solo se dispone de una cota inferior y de la mejor
solución encontrada hasta ahora. De las 242 instancias, 156 tienen óptimo
conocido (su cota inferior publicada es igual a su mejor solución conocida).
Esta diversidad es deseable para un portafolio, porque el desempeño relativo
de los \textit{solvers} suele cambiar con el tamaño y la estructura de la
instancia.

\subsection{Conversión al formato de MiniZinc}

Las instancias se encuentran en el formato de texto de la OR-Library, compuesto
por una primera línea con $n$ y $m$ seguida de una línea por trabajo con $m$
pares \textit{(máquina, duración)} en orden, con las máquinas numeradas desde
cero. El programa \texttt{src/01\_convertir\_instancias.py} las transforma en
archivos \texttt{.dzn} con los parámetros del modelo
(Sección~\ref{sec:modelo_original}), sin modificar los índices, y genera una
segunda versión con los parámetros \texttt{lb} y \texttt{ub}
(Sección~\ref{sec:cotas_instancias}). Se verificó que las matrices de las 242
instancias tuvieran las dimensiones declaradas y que los índices de máquina
estuvieran en el rango $0,\dots,m-1$, y, como prueba de extremo a extremo,
todas las configuraciones resolvieron \texttt{ft06} y obtuvieron su óptimo
conocido, 55.

%------------------------------------------------------------------------------
\section{Cotas de las instancias}
\label{sec:cotas_instancias}
%------------------------------------------------------------------------------

Para cada instancia se calcularon la cota inferior $lb$ y la cota superior
$ub$ de la Sección~\ref{sec:cotas} con el programa
\texttt{src/02\_generar\_cotas.py}. La calidad alcanzada en cada familia frente
a la mejor solución conocida (BKS, Sección~\ref{sec:bks_gap}) se reporta en la
Tabla~\ref{tab:cotas}.

\begin{table}[H]
  \centering
  \small
  \singlespacing
  \caption{Calidad de las cotas por familia.}
  \label{tab:cotas}
  \begin{tabular}{lrrrrrl}
    \toprule
    Familia & Inst. & lb/BKS (\%) & Gap de ub (\%) & $H$/BKS & lb $=$ ub & Regla más frecuente \\
    \midrule
    ft & 3 & 84.0 & 23.1 & 4.5 & 0 & MWKR \\
    la & 40 & 93.5 & 14.7 & 6.6 & 4 & MWKR \\
    abz & 5 & 80.4 & 23.1 & 9.2 & 0 & FIFO \\
    orb & 10 & 70.3 & 26.5 & 5.6 & 0 & FIFO \\
    swv & 20 & 88.1 & 32.4 & 8.3 & 2 & FIFO \\
    yn & 4 & 77.0 & 25.4 & 13.0 & 0 & FIFO \\
    ta & 80 & 90.8 & 22.8 & 13.9 & 0 & FIFO \\
    dmu & 80 & 86.7 & 34.4 & 13.1 & 0 & MWKR \\
    \midrule
    Total & 242 & 88.3 & 26.3 & 11.4 & 6 & MWKR \\
    \bottomrule
  \end{tabular}
  \par\vspace{2pt}{\footnotesize lb: cota inferior; ub: mejor programa de Giffler-Thompson; $H$: horizonte del modelo original (suma de duraciones).}
\end{table}


Las cotas reducen el espacio de búsqueda en un orden de magnitud, puesto que el
horizonte trivial del modelo original es, en promedio, 11.4 veces la mejor
solución conocida, mientras que la cota superior la supera solo en un
26.3\,\%. El efecto es mayor en las familias de instancias grandes
(\texttt{ta}, \texttt{dmu}, \texttt{yn}), donde el horizonte trivial es más de
13 veces la BKS. La heurística ya es, además, una solución razonable, de modo
que el portafolio parte en el peor caso de un programa factible a una
distancia moderada de la mejor solución conocida. Las seis instancias en que
$lb = ub$ no requieren ejecutar ningún \textit{solver} y se excluyen de la
comparación entre estrategias, que se hace sobre las 236 restantes. El mejor
resultado se reparte entre varias reglas de prioridad, ya que MWKR, FIFO y MOR
obtienen la menor cota en distintas familias, lo que justifica evaluar las seis
reglas por instancia.

%------------------------------------------------------------------------------
\section{Extracción de características}
\label{sec:features}
%------------------------------------------------------------------------------

Cada instancia se describe mediante las características genéricas que calcula
\texttt{mzn2feat} a partir del modelo MiniZinc y un grupo de características
propias del JSSP calculadas directamente sobre los datos de la instancia.

\subsection{Características genéricas con \texttt{mzn2feat}}

Se usó \texttt{mzn2feat-1.0} \citep{amadini2014mzn2feat, amadini2015sunnycp},
el extractor de
características del portafolio SUNNY-CP, que traduce el modelo y los datos a
FlatZinc y calcula 95 características numéricas sobre el resultado. Su uso en
portafolios de programación por restricciones está documentado, lo que hace
comparables los resultados con los de SUNNY, y sus características se calculan
exclusivamente a partir de la estructura del modelo que recibe el
\textit{solver}. Como todas las instancias usan el mismo modelo, las
diferencias entre sus vectores se deben solo a los datos. El programa
\texttt{src/03\_caracteristicas\_mzn2feat.py} ejecutó la extracción sobre el
modelo original dentro del contenedor Docker oficial de SUNNY-CP
(\texttt{jacopomauro/sunny-cp}), que incluye la herramienta compilada y
MiniZinc~1.6, ya que \texttt{mzn2feat} se distribuye para versiones antiguas
de MiniZinc. Se usó el modelo original porque el propósito es describir la
instancia, y las cotas se derivan de esos mismos datos. Las características se
agrupan en seis categorías, identificadas por el prefijo de su nombre
(Tabla~\ref{tab:categorias_features}), y dentro de cada una se calculan
estadísticos como la media (\texttt{avg}), el coeficiente de variación
(\texttt{cv}) o la entropía (\texttt{ent}) del grado o del tamaño de dominio
de variables y restricciones (Anexo~\ref{anx:caracteristicas}).

\begin{table}[H]
  \centering
  \small
  \singlespacing
  \caption{Categorías de características de \texttt{mzn2feat} y número de
  características activas tras el preprocesamiento
  (Sección~\ref{sec:limpieza}).}
  \label{tab:categorias_features}
  \begin{tabularx}{\textwidth}{>{\ttfamily}l >{\raggedright\arraybackslash}X c c}
    \toprule
    \normalfont\textbf{Prefijo} & \textbf{Describe} & \textbf{Extraídas} & \textbf{Activas} \\
    \midrule
    v\_  & Variables: número, grado (en cuántas restricciones aparece cada una) y tamaño de dominio & 27 & 26 \\
    d\_  & Tipos de dominio: proporción de variables y restricciones enteras, booleanas, de conjunto y de punto flotante & 18 & 8 \\
    c\_  & Restricciones: número, grado y tamaño de los dominios de sus variables & 27 & 21 \\
    gc\_ & Restricciones globales: cuántas hay y de qué tipo & 4 & 3 \\
    o\_  & Función objetivo: dominio y grado de la variable a optimizar & 8 & 7 \\
    s\_  & Anotaciones de búsqueda incluidas en el modelo & 11 & 0 \\
    \midrule
    \normalfont\textbf{Total} & & 95 & 65 \\
    \bottomrule
  \end{tabularx}
\end{table}

Las 11 características de búsqueda (\texttt{s\_}) son constantes, porque el
modelo original no incluye anotaciones de búsqueda, y lo mismo ocurre con las
que miden tipos de variables o restricciones que el modelo no usa. Al
emplearse un único modelo, muchas de las restantes dependen sobre todo del
tamaño de la instancia. Por ejemplo, el dominio de las variables de inicio es
$0..H$, de modo que los estadísticos de tamaño de dominio crecen con el número
de operaciones y con sus duraciones. Pocas describen la estructura propia del
JSSP, como la distribución de la carga entre máquinas o el orden de las rutas.

\subsection{Características propias del JSSP}
\label{sec:features_propias}

Otro aporte de este trabajo es el diseño de un conjunto de 20 características
propias del JSSP para complementar a \texttt{mzn2feat}, las cuales se calculan
directamente sobre las matrices de máquinas y duraciones de cada instancia con
el programa \texttt{src/\allowbreak 04\_\allowbreak caracteristicas\_\allowbreak propias.py}.
Siguiendo la línea de los trabajos que caracterizan instancias del JSSP a
partir de su estructura \citep{mirshekarian2016}, estos descriptores capturan
el tamaño, la distribución de las duraciones, el equilibrio de la carga entre
máquinas y trabajos, y la forma de las rutas
(Tabla~\ref{tab:features_propias}). Su cálculo toma menos de un segundo para
las 242 instancias y no requiere MiniZinc.

\begin{table}[H]
  \centering
  \small
  \singlespacing
  \caption{Características propias del JSSP. $P$: duraciones; $L_k$: carga
  (suma de duraciones) de la máquina $k$; $W_j$: duración total del trabajo
  $j$; $LB$: cota inferior.}
  \label{tab:features_propias}
  \begin{tabularx}{\textwidth}{l >{\ttfamily\raggedright\arraybackslash}p{4.6cm} >{\raggedright\arraybackslash}X}
    \toprule
    \textbf{Grupo} & \normalfont\textbf{Característica} & \textbf{Definición} \\
    \midrule
    Tamaño & jf\_n, jf\_m & Número de trabajos y de máquinas \\
     & jf\_ratio\_nm, jf\_ops & $n/m$ y número de operaciones $n \cdot m$ \\
    \midrule
    Duraciones & jf\_p\_mean & Duración media de las operaciones \\
     & jf\_p\_cv & Coeficiente de variación de las duraciones \\
     & jf\_p\_range & $(\max P - \min P) / \overline{P}$ \\
    \midrule
    Carga & jf\_load\_cv & Coeficiente de variación de $L_k$ \\
     & jf\_bottleneck & $\max L_k / \overline{L}$: intensidad del cuello de botella \\
     & jf\_job\_cv & Coeficiente de variación de $W_j$ \\
     & jf\_jobmax\_rel & $\max W_j / \overline{W}$ \\
     & jf\_machine\_vs\_job & $\max L_k / \max W_j$: qué componente domina la cota inferior \\
     & jf\_lb\_over\_meanload & $LB / \overline{L}$ \\
     & jf\_lb\_over\_horizon & $LB / H$: holgura del horizonte trivial \\
     & jf\_mach\_pmean\_cv & Coeficiente de variación de la duración media por máquina \\
    \midrule
    Rutas & jf\_route\_pos\_entropy & Entropía media de las máquinas que ocupan cada posición de las rutas \\
     & jf\_route\_kendall & Concordancia media (tipo $\tau$ de Kendall) entre los órdenes de máquinas de dos trabajos; vale 1 en un \textit{flow shop} \\
     & jf\_first\_mach\_ent, jf\_last\_mach\_ent & Entropía de la primera y de la última máquina de las rutas \\
     & jf\_bneck\_pos & Posición media, en las rutas, de las operaciones de la máquina más cargada \\
    \bottomrule
  \end{tabularx}
\end{table}

En la validación del Capítulo~\ref{cap:resultados} se comparan tres conjuntos
conformados por las 65 características activas de \texttt{mzn2feat}, las 20
propias y la unión de ambos (85 características). La última característica de
la tabla resultó redundante, ya que cuando cada máquina procesa exactamente $n$
operaciones, como ocurre en todas las instancias clásicas, coincide con
\texttt{jf\_load\_cv}. Se conservó para mantener el conjunto tal como se
definió antes de la evaluación, con el único efecto de dar un peso doble a esa
magnitud en la distancia de los métodos de vecinos más cercanos.

%------------------------------------------------------------------------------
\section{Ejecución de los solvers}
\label{sec:ejecucion_solvers}
%------------------------------------------------------------------------------

Se realizaron dos campañas de ejecución de 300 segundos por instancia y
configuración. La primera, con el modelo original, sirvió como diagnóstico y
motivó la incorporación de las cotas. La segunda, con los modelos con cotas,
es la que alimenta el portafolio. En ambas campañas se emplearon tres
\textit{solvers} a través de MiniZinc~2.9.1, cuyas versiones y técnicas
principales se indican en la Tabla~\ref{tab:solvers}.

\begin{table}[H]
  \centering
  \small
  \singlespacing
  \caption{\textit{Solvers} utilizados.}
  \label{tab:solvers}
  \begin{tabularx}{\textwidth}{l c >{\raggedright\arraybackslash}X}
    \toprule
    \textbf{\textit{Solver}} & \textbf{Versión} & \textbf{Técnica principal} \\
    \midrule
    Chuffed & 0.13.2 & Generación perezosa de cláusulas (\textit{Lazy Clause Generation}) \\
    Gecode  & 6.3.0  & Propagación y búsqueda sobre dominios finitos \\
    OR-Tools CP-SAT & 9.12.4544 & Híbrido: generación perezosa de cláusulas sobre un núcleo SAT, combinada con técnicas de programación lineal \\
    \bottomrule
  \end{tabularx}
\end{table}

\subsection{Protocolo común}

Todas las ejecuciones siguieron el mismo protocolo. MiniZinc impone el límite
de 300 segundos con la opción \texttt{--time-limit}, que incluye la traducción a
FlatZinc y la búsqueda, y el programa de ejecución termina el árbol de procesos
si se supera ese límite en más de 20 segundos. Se usó la opción \texttt{-a},
para que el \textit{solver} reporte cada solución mejorada (sin ella, una
ejecución que agota el tiempo no reporta ninguna), y \texttt{-f}, que permite a
cada \textit{solver} usar su propia estrategia, salvo en el modelo con
estrategia de búsqueda. Ningún \textit{solver} se ejecutó con varios hilos, lo
que se verificó con la razón entre el tiempo de CPU y el de pared en la
instancia \texttt{dmu01} (0.92 para Chuffed, 0.93 para Gecode y 0.92 para
CP-SAT). Para cada par (instancia, configuración) se registra el último
\textit{makespan} reportado, el tiempo de pared, si se probó la optimalidad y
el estado final (\texttt{OPTIMAL}, \texttt{SATISFIABLE} si se encontró al
menos una solución sin probar que es óptima, o \texttt{UNKNOWN} si no se
encontró ninguna).

\subsection{Campaña con el modelo original}

La primera campaña ejecutó Chuffed, Gecode y CP-SAT con el modelo original
sobre las 242 instancias. Los resultados globales de esta exploración inicial
se sintetizan en la Tabla~\ref{tab:desempeno_300}, donde el \textit{gap}
corresponde a la definición de la Sección~\ref{sec:bks_gap}.

\begin{table}[H]
  \centering
  \small
  \singlespacing
  \caption{Desempeño de los \textit{solvers} con el modelo original
  (242 instancias, 300 segundos).}
  \label{tab:desempeno_300}
  \begin{tabular}{l c c c}
    \toprule
    & \textbf{Chuffed} & \textbf{Gecode} & \textbf{CP-SAT} \\
    \midrule
    Instancias con solución             & 242     & 242    & 242 \\
    Óptimos probados                     & 32      & 7      & 69 \\
    Instancias con \textit{gap} de 0\,\% & 32      & 7      & 76 \\
    \textit{Gap} medio (\%)              & 637.06  & 25.03  & 4.22 \\
    \textit{Gap} mediano (\%)            & 580.30  & 26.06  & 2.79 \\
    Tiempo medio (s)                     & 269.80  & 293.08 & 228.31 \\
    \bottomrule
  \end{tabular}
  \par\vspace{2pt}{\footnotesize Sin las seis instancias con $lb = ub$, el
  \textit{gap} medio es de 648.3, 25.6 y 4.3\,\% (Tabla~\ref{tab:efecto_cotas}).}
\end{table}

CP-SAT obtiene un \textit{makespan} igual o mejor que el de Chuffed y Gecode en
las 242 instancias, y estrictamente mejor en 210. Su ventaja generalizada sobre
ambos hace innecesaria una regla de selección, por lo que se dejó como
referencia externa (Sección~\ref{sec:composicion}). Chuffed prueba más óptimos
que Gecode en las instancias pequeñas, pero en 199 instancias su \textit{gap}
supera el 100\,\%, dado que con el horizonte trivial acepta como primeras
soluciones programas casi secuenciales y no logra mejorarlos dentro del límite.
Este comportamiento obedece a la amplitud del dominio en la formulación inicial
y motivó la incorporación de las cotas (Sección~\ref{sec:cotas}). Con ese
modelo el portafolio tenía también poco margen, ya que Gecode obtiene el mejor
\textit{makespan} entre los dos en 206 instancias, Chuffed en 29 y empatan en 7,
y la ventaja de Chuffed se limita a instancias pequeñas en las que ambos están
cerca del óptimo.
Estos \textit{solvers}, con el modelo original y su búsqueda por defecto, son
el punto de comparación de la propuesta, y la Tabla~\ref{tab:portafolios} los
contrasta con el portafolio construido sobre el modelo con cotas.

\subsection{Campaña con los modelos con cotas}

La segunda campaña ejecutó las tres configuraciones del portafolio
(Sección~\ref{sec:composicion}), integradas por Chuffed y Gecode con el modelo
con cotas y por Chuffed con el modelo con estrategia de búsqueda (Chuffed V2).
Se usó el programa de ejecución \texttt{src/05\_ejecutor}, diseñado para que los
resultados fueran verificables y la campaña pudiera repetirse. Antes de
empezar, comprueba las versiones de MiniZinc y de los \textit{solvers}
(Tabla~\ref{tab:solvers}) y resuelve \texttt{ft06}, cuyo óptimo (55) es
conocido. Con la opción \texttt{--output-time} de MiniZinc registra la
\textit{trayectoria} de cada ejecución, es decir, el instante en que se
encontró cada solución mejorada, que permite conocer $C(a, i, \tau)$ para
cualquier $\tau$ y evaluar cualquier \textit{schedule} sin volver a ejecutar
los \textit{solvers} (Sección~\ref{sec:formulacion_schedule}). También
comprueba que cada \textit{makespan} esté entre $lb$ y $ub$, marca las
ejecuciones sin solución y guarda cada resultado en cuanto termina, de modo
que una campaña interrumpida continúa donde quedó.

Las 708 ejecuciones (236 instancias por tres configuraciones) terminaron sin
errores. Ningún \textit{makespan} quedó fuera del intervalo $[lb, ub]$, y
todos los óptimos declarados coinciden con la mejor solución conocida. Una
ejecución sin solución significa que la configuración no encontró, en 300
segundos, un programa al menos tan bueno como el heurístico, situación que
ocurrió en 12 instancias con Chuffed, todas de 750 operaciones o más, en 24 con
Chuffed V2 y en 65 con Gecode, en ambos casos sobre todo en instancias grandes.
En esos casos el resultado de la configuración es la $ub$. Los resultados de
esta campaña se analizan en el Capítulo~\ref{cap:resultados} y se listan por
instancia en el Anexo~\ref{anx:resultados}.

Con el mismo programa y el mismo protocolo se ejecutó también CP-SAT con el
modelo con cotas, como \textit{solver} de referencia. Sus 236 ejecuciones
fueron válidas y encontraron solución en todas las instancias.

%------------------------------------------------------------------------------
\section{Preprocesamiento y base de conocimiento}
\label{sec:preprocesamiento}
%------------------------------------------------------------------------------

El preprocesamiento cruza por nombre de instancia las cotas, las mejores
soluciones conocidas, las características y los resultados de la campaña con
cotas. Las 242 instancias se vincularon en todas ellas.

\subsection{Mejores soluciones conocidas y brecha de optimalidad}
\label{sec:bks_gap}

Para medir la calidad de una solución se compara su \textit{makespan} con la
mejor solución conocida (BKS, por \textit{Best Known Solution}) de la
literatura, tomada del repositorio de \citet{weise2019jssp}, que recopila los
mejores resultados publicados hasta 2020. Para cada instancia $i$ y
configuración $a$, la brecha de optimalidad o \textit{gap} es
\begin{equation}
  \mathit{gap}(a,i) = 100 \cdot \frac{C_{\max}(a,i) - \mathit{BKS}(i)}{\mathit{BKS}(i)},
  \label{eq:gap}
\end{equation}
expresado en porcentaje, donde $C_{\max}(a,i)$ es el \textit{makespan}
obtenido por $a$ en $i$. Un \textit{gap} de 0\,\% indica que la configuración
alcanzó la mejor solución conocida. En las 156 instancias con óptimo
conocido, la BKS coincide con el óptimo, mientras que en las demás corresponde
a la mejor solución publicada, de modo que el \textit{gap} es una cota inferior
de la distancia real al óptimo.

\subsection{Limpieza y normalización de características}
\label{sec:limpieza}

\texttt{mzn2feat} reporta como \texttt{nan} o \texttt{inf} algunos
estadísticos que no se pueden calcular, por ejemplo cuando implican una
división por cero. Estos valores se trataron como faltantes. Solo los tenían
\texttt{c\_cv\_domdeg\_cons}, en 117 instancias, y \texttt{o\_dom\_std}, en
241, y se imputan con la mediana de la característica en el conjunto de
entrenamiento. Se eliminaron después las 30 características de
\texttt{mzn2feat} con varianza cero en las 242 instancias (anotaciones de
búsqueda, tipos de variables y restricciones que el modelo no usa y algunos
estadísticos que el modelo fija), con lo que quedaron 65 activas
(Tabla~\ref{tab:categorias_features}). Las 20 características propias no
tienen valores faltantes ni constantes.

Por último, siguiendo a SUNNY \citep{amadini2014sunny}, cada característica
$x$ se escala al intervalo $[-1, 1]$,
\begin{equation}
  x' = 2\,\frac{x - x_{\min}}{x_{\max} - x_{\min}} - 1,
  \label{eq:normalizacion}
\end{equation}
donde $x_{\min}$ y $x_{\max}$ son el mínimo y el máximo en el conjunto de
entrenamiento, y los valores de una instancia nueva que quedan fuera de ese
rango se recortan a $[-1, 1]$. Así se evita que las características con
valores grandes dominen la distancia euclidiana de los métodos de vecinos
más cercanos. La imputación y la normalización se calculan dentro de cada
partición de entrenamiento, para que la instancia evaluada no influya en su
propia transformación (Sección~\ref{sec:validacion}).

\subsection{Base de conocimiento}

Un aporte metodológico de este trabajo es la construcción de una base de
conocimiento con las trayectorias completas de soluciones de cada
configuración, lo que permite evaluar de forma exacta cualquier
\textit{schedule} por simulación sin repetir la ejecución de los
\textit{solvers}. Para cada una de las 242 instancias, la base reúne sus datos
generales (familia, $n$, $m$, número de operaciones, $lb$, $ub$, regla que
produjo la $ub$ y BKS), sus 85 características (65 de \texttt{mzn2feat} y 20
propias) y, para cada configuración, el \textit{makespan} final, si probó el
óptimo y en qué tiempo, el estado de la ejecución y la trayectoria completa
de soluciones. A diferencia de la base de conocimiento de SUNNY, que guarda
solo el resultado final de cada \textit{solver}, las trayectorias indican
cuánto tiempo necesitó cada configuración para alcanzar cada nivel de calidad,
que es justamente la información que requiere un \textit{schedule}. La base de
conocimiento se construye con el programa \texttt{src/06\_analisis\_campana.py},
que además produce los análisis descriptivos del
Capítulo~\ref{cap:resultados}. Con esto concluye la construcción del conjunto
de datos de entrenamiento (actividades A8 a A10).
